\subsection{BÀI TẬP TRẮC NGHIỆM}

\ind{PHẦN I.} \inden{Câu trắc nghiệm nhiều phương án lựa chọn. Mỗi câu hỏi học sinh chỉ chọn một phương án.}\\
\setcounter{ex}{0}
\Opensolutionfile{ans}[ans/1D6-B2-1]
\begin{ex}
	Cho $a$ là số thực dương khác $1$. Khẳng định nào dưới đây là \textbf{sai}?
	\choice
	{ $\log_a {2}\cdot \log_2 {a}=1$}
	{$\log_a {1}=0$}
	{$\log_a {a}=1$}
	{\True $\log_a {2}=\dfrac{1}{\log_a {2}}$}
	
	\loigiai{$\log_a 2=\dfrac{1}{\log_a {2}}$ không phải là công thức đổi cơ số, nếu đúng là $\log_a {2}=\dfrac{1}{\log_2 {a}}$. }
	
\end{ex} 

\begin{ex}
	Cho các số thực $a, b>1$. Mệnh đề nào sau đây là mệnh đề đúng?
	\choice
	{$\log_a\dfrac{a}{b}=\log_b a$}
	{$\log_a\dfrac{a}{b}=1+\log_a b$}
	{$\log_a\dfrac{a}{b}=\log_a b$}
	{\True $\log_a\dfrac{a}{b}=1-\log_a b$}
	\loigiai{Ta có $\log_a\dfrac{a}{b}=\log_a a-\log_a b=1-\log_a b$.
	}
\end{ex}

\begin{ex}
	Cho $a$ là số thực dương tùy ý khác $ 1 $. Mệnh đề nào dưới đây đúng?
	\choice
	{$\log_{2}a=\log_{a}2$}
	{$\log_{2}a=\dfrac{1}{\log_{2}a}$}
	{\True $\log_{2}a=\dfrac{1}{\log_{a}2}$}
	{$\log_{2}a=-\log_{a}2$}
	\loigiai{
		Với $a$ là số thực dương tùy ý khác $ 1 $, mệnh đề đúng là mệnh đề: $\log_{2}a=\dfrac{1}{\log_{a}2}$ 
	}
\end{ex}

\begin{ex}
	Với $a,b,c$ là các số thực dương khác $1$, mệnh đề nào dưới đây là mệnh đề \textbf{sai}?
	\choice
	{$\log_ab=\dfrac{\log b}{\log a}$}
	{\True $\log_ab=\dfrac{\log_ca}{\log_cb}$}
	{$\log_ab=\dfrac{1}{\log_ba}$}
	{$\log_ab=\dfrac{\ln b}{\ln a}$}
	\loigiai{
		Với $a,b,c$ dương và khác $1$, $\log_ab=\dfrac{\log_cb}{\log_ca}$ mới là công thức đổi cơ số đúng.}
\end{ex}

\begin{ex}
	Cho $a,b>0$. Tìm mệnh dề đúng trong các mệnh đề sau.
	\choice
	{\True $\ln \dfrac{a}{b}=\ln a+\ln \dfrac{1}{b}$}
	{$\ln \dfrac{a}{b}=\ln b-\ln a$}
	{$\ln \dfrac{a}{b}=\dfrac{\ln a}{\ln b}$}
	{$\ln \dfrac{a}{b}=\ln a-\ln \dfrac{1}{b}$}
	\loigiai{
		Ta có $\ln \dfrac{a}{b}=\ln \left(a\cdot \dfrac{1}{b}\right)=\ln a+\ln \dfrac{1}{b}.$
	}
\end{ex} 

\begin{ex}
	Giá trị của biểu thức $A=4^{\log_27}$ bằng
	\choice
	{$14$}
	{$28$}
	{$2$}
	{\True $49$}
	\loigiai{ Ta có: $A=4^{\log_27}=2^{2\log_27}=2^{\log_27^2}=49$.
	}
\end{ex}

\begin{ex}
	Biết $\log_6 a =2$. Tính $I=\log_a 6$.
	\choice
	{$I=36$}
	{\True $I=\dfrac{1}{2}$}
	{$I=64$}
	{$I=\dfrac{1}{4}$}
	\loigiai{ Ta có $I=\log_a 6 =\dfrac{1}{\log_6 a}=\dfrac{1}{2}$. }	
	
\end{ex}

\begin{ex}
	Tính giá trị của biểu thức $I = a \cdot \log_2 \sqrt{8}$.
	\choice
	{$I = \dfrac{2}{3}$}
	{\True $I = \dfrac{3a}{2}$}
	{$I = \dfrac{2a}{3}$}
	{$I = \dfrac{3}{2}$}
	\loigiai
	{
		$I = a \cdot \log_2 \sqrt{8} = a \cdot \log _2 2^{\frac{3}{2}} = \dfrac{3}{2} \cdot a \cdot \log_2 2 = \dfrac{3a}{2}$.
	}
\end{ex}

\begin{ex}
	Tính giá trị của biểu thức $N=\log_a\sqrt{a\sqrt{a}}$ với $0<a\ne 1$.
	\choice
	{$N=\dfrac{-3}{4}$}
	{$N=\dfrac{4}{3}$}
	{$N=\dfrac{3}{2}$}
	{\True $N=\dfrac{3}{4}$}
	\loigiai{
		Ta có: $N=\log_a\sqrt{a\sqrt{a}}=\log_a\sqrt{a^{\frac{3}{2}}}=\log_aa^{\frac{3}{4}}=\dfrac{3}{4}$.}
\end{ex}

\begin{ex}
	Biểu thức $ \log_2\left(2\sin\dfrac{\pi}{12}\right)+ \log_2\left(\cos\dfrac{\pi}{12}\right)$ có giá trị bằng
	\choice
	{$ -2 $}
	{\True $ -1 $}
	{$ 1 $}
	{$ \log_2\sqrt{3}-1 $}
	\loigiai{
		Ta có: $ \log_2\left(2\sin\dfrac{\pi}{12}\right)+ \log_2\left(\cos\dfrac{\pi}{12}\right)=\log_2\left(2\sin\dfrac{\pi}{12}\cdot\cos\dfrac{\pi}{12}\right)=\log_2\left(\sin\dfrac{\pi}{6}\right)=\log_2\dfrac{1}{2}=-1.$
	}
\end{ex}

\begin{ex}
	Biết $\log_2x=a$, tính theo $a$ giá trị của biểu thức $P=\log_2\left( 4x^2\right) $.
	\choice
	{$P=2+a$}
	{$P=4+2a$}
	{$P=4+a$}
	{\True $P=2+2a$}
	\loigiai{Có $P=\log_24x^2=\log_24+\log_2x^2=2+2\log_2x=2+2a$.}
\end{ex}


\begin{ex}
	Cho $\log_c a=2$ và $\log_c b=4$. Tính $P=\log_a b^4$.
	\choice
	{\True $P=8$}
	{$P=\dfrac{1}{32}$}
	{$P=\dfrac{1}{8}$}
	{$P=32$}
	\loigiai{
		Ta có $\dfrac{\log_c a}{\log_c b}=\log_a b=\dfrac{4}{2}=2 \Rightarrow 4\log_a b=\log_a b^4=8$. 
	}
\end{ex}

\begin{ex}
	Cho $\log _a b = 5$, $\log_a c = -3$. Giá trị biểu thức $\log_a \left( \dfrac{a^4 \sqrt[3]{b}}{c^2} \right)$ là
	\choice
	{$-\dfrac{1}{3}$}
	{$-40$}
	{$40$}
	{\True $\dfrac{35}{3}$}
	\loigiai{
		Ta có $\log_ab = 5 \Rightarrow b = a^5$, $\log_ac = -3 \Rightarrow c = a^{-3}$.\\
		Vậy $\log_a \left( \dfrac{a^4 \sqrt[3]{b}}{c^2} \right) = \log_a\left(\dfrac{a^4\cdot a^{\frac{5}{3}}}{a^{-6}}\right) = \log_a a^{\frac{35}{3}} = \dfrac{35}{3}$.
	}
\end{ex}

\begin{ex} (TN -- 2021)
	Với mọi $a$, $b$ thỏa mãn $\log_{2}a^{3}+\log_{2}b=6$. Khẳng định nào dưới đây đúng?
	\choice
	{$a^{3}b=64$}
	{$a^{3}b=36$}
	{\True $a^{3}+b=64$}
	{$a^{3}+b=36$}
	\loigiai{
		Ta có: $\log_{2}a^{3}+\log_{2}b=\log_{2}\left(a^{3}b\right)=6$
		$\Rightarrow a^{3}b=2^{6}=64$.}
\end{ex}

\begin{ex}
	Nếu $a=\log_{30}3$ và $b=\log_{30}5$ thì
	\choice
	{$\log_{30}1350=a+2b+1$}
	{\True $\log_{30}1350=2a+b+1$}
	{$\log_{30}1350=a+2b+2$}
	{$\log_{30}1350=2a+b+2$}
	\loigiai{
		Ta có:$\log_{30}1350=\log_{30}\left(30.3^{2}.5\right)=1+\log_{30}3^{2}+\log_{30}5=1+2\log_{30}3+\log_{30}5=1+2a+b$
	}
\end{ex}

\begin{ex}
	Cho $\log_2{7}=a$, $\log_3{7}=b$. Tính $\log_6{7}$ theo $a$ và $b$.
	\choice
	{$\dfrac{1}{a+b}$}
	{$a^2+b^2$}
	{$a+b$}
	{\True $\dfrac{ab}{a+b}$}
	\loigiai{
		$\log_67=\dfrac{1}{\log_76}=\dfrac{1}{\log_72+\log_73}=\dfrac{1}{\dfrac{1}{a}+\dfrac{1}{b}}=\dfrac{ab}{a+b}$.
	}
\end{ex}

\begin{ex}
	Cho $a=\log _3 15$, $b=\log _3 10$. Tính $\log _{\sqrt{3}} 50$ theo $a$ và $b$.
	\choice
	{\True $\log _{\sqrt{3}} 50=2\left( {a+b-1} \right)$}
	{$\log _{\sqrt{3}} 50=4\left( {a+b+1} \right)$}
	{$\log _{\sqrt{3}} 50=a+b-1$}
	{$\log _{\sqrt{3}} 50=3\left( {a+b+1} \right)$}
	\loigiai{
		Ta có $a=\log _315=1+\log _35 \Rightarrow \log _35=a-1$.
		\\Vậy $\log _{\sqrt{3}}50=2 \log _350=2\left( {\log _3\left( {5 \cdot 10} \right)} \right)=2\left( {\log _35+\log _310} \right)=2\left( {a+b-1} \right)$.}
\end{ex}

\begin{ex}
	Cho $\log_2 6 =a; \log_2 7 =b$. Tính $\log_3 7$ theo $a$ và $b$.
	\choice
	{\True $\log_3 7 = \dfrac{b}{a-1}$}
	{$\log_3 7 = \dfrac{a}{b-1}$}
	{$\log_3 7 = \dfrac{b}{1-a}$}
	{$\log_3 7 = \dfrac{a}{1-b}$}
	\loigiai{
		$ \log_3 7 = \dfrac{\log_2 7}{\log_2 3} = \dfrac{\log_2 7}{\log_2 6 - 1} = \dfrac{b}{a-1}$.
	}
\end{ex}

\begin{ex}
	Đặt $a=\log_{12}6$, $b=\log_{12} 7$. Hãy biểu diễn $\log_{2} 7$ theo $a$ và $b$.
	\choice
	{$\dfrac{b}{a+1}$}
	{\True $\dfrac{b}{1-a}$}
	{$\dfrac{a}{b-1}$}
	{$\dfrac{a}{b+1}$}
	\loigiai{
		$\log_2 7=\dfrac{\log_{12}7}{\log_{12}2}=\dfrac{\log_{12}7}{\log_{12}\dfrac{12}{6}}=\dfrac{\log_{12}7}{1-\log_{12}6}=\dfrac{b}{1-a}$.
	}
\end{ex}

\begin{ex}
	Đặt $a=\ln 2; b=\ln 5$. Hãy biểu diễn $I=\ln\dfrac{1}{2}+\ln\dfrac{2}{3}+...+\ln\dfrac{98}{99}+\ln\dfrac{99}{100}$ theo $a$ và $b$.
	\choice
	{\True $I=-2(a+b)$}
	{$I=2(a+b)$}
	{$I=-2(a-b)$}
	{$I=2(a-b)$}
	\loigiai{Ta có $I=\ln\dfrac{1}{2}+\ln\dfrac{2}{3}+...+\ln\dfrac{98}{99}+\ln\dfrac{99}{100}=\ln\dfrac{1}{100}=-\ln 100=-2\ln(2.5)=-2(\ln 2+\ln 5)=-2(a+b)$}
\end{ex}

\Closesolutionfile{ans}

\ind{PHẦN II.} \inden{Câu trắc nghiệm đúng sai. Trong mỗi ý a), b), c), d) ở mỗi câu, học sinh chọn đúng hoặc sai.}\\
\Opensolutionfile{ans}[ans/1D6-B2-2]

\begin{ex}
	Cho biểu thức $B=2 \ln \sqrt{e x}-\ln \frac{e^2}{\sqrt{x}}+\ln 3 \cdot \log _3\left(e x^2\right)$.
	\choiceTF
	{\True Nếu $\ln x=2$ thì $B=7$}
	{\True Nếu $\ln x=4$ thì $B=14$}
	{\True Nếu $\ln x=3$ thì $B=\dfrac{15}{2}$}
	{Nếu $\ln x=6$ thì $B=18$}
	\loigiai{
		
	}
\end{ex}

\begin{ex}%[Câu 1]%[1D6H2-2]
	Cho các biểu thức sau $P=\log _2 8+\log _3 27-\log _5 5^3$; $Q=\ln (2\mathrm{e})-\log 100$. Xét tính đúng sai của các khẳng định sau:
	\choiceTF
	{$P+Q=2\ln 2$}
	{\True $Q-P=\ln 2-4$}
	{\True $3Q+P=3\ln 2$}
	{\True $2Q+P=2\ln 2+1$}
	\loigiai{
		Ta có $P=\log _2 8+\log _3 27-\log _5 5^3=\log _2 2^3+\log _3 3^3-\log _5 5^3=3+3-3=3$.\\
		Ta có $Q=\ln (2\mathrm{e})-\log 100=\ln 2+\ln \mathrm{e}-\log{10^2}=\ln 2+1-2=\ln 2-1$.
		\begin{itemchoice}
			\itemch Sai. Vì $P+Q=3+(\ln 2-1)=2+\ln 2$.
			\itemch Đúng. Vì $Q-P=(\ln 2-1)-3=\ln 2-4$.
			\itemch Đúng. Vì $3Q+P=3(\ln 2-1)+3=3\ln 2$.
			\itemch Đúng. Vì $2Q+P=2(\ln 2-1)+3=2\ln 2+1$.
		\end{itemchoice}
	}
\end{ex}
\begin{ex}%[Câu 2]%[1D6H2-2]
	Cho các biểu thức sau $A=\log_{2^{2030}}4-\dfrac{1}{1015}+\ln{\mathrm{e}^{2035}}$; $B=\log_53\cdot{\log_2}5-\dfrac{\ln 9}{\ln 4}$. Xét tính đúng sai của các khẳng định sau:
	\choiceTF
	{\True $A$ chia hết cho $5$}
	{$A-B=2036$}
	{\True $A+2024B=2035$}
	{\True $A-2024B=2035$}
	\loigiai{
		Ta có
		\begin{eqnarray*}
			&A&=\log_{2^{2030}}4-\dfrac{1}{1015}+\ln{\mathrm{e}^{2035}}\\
			&&=\log_{2^{2030}}2^2-\dfrac{1}{1015}+2035\\
			&&=\dfrac{2}{2030}-\dfrac{1}{1015}+2035=2035.
		\end{eqnarray*}
		Ta có
		\begin{eqnarray*}
			&B&=\log_53.\log_25-\dfrac{\ln 9}{\ln 4}\\
			&&=\log_25.\log_53-\log_49\\
			&&=\log _2 3-\log _{2^2} 3^2\\
			&&=\log _2 3-\log _2 3=0.
		\end{eqnarray*}
		Do đó
		\begin{itemchoice}
			\itemch Đúng. Vì $A=2035$ chia hết cho $5$.
			\itemch Sai. Vì $A-B=2035-0=2035$.
			\itemch Đúng. Vì $A+2024B=2035+2024\cdot 0=2035$.
			\itemch Đúng. Vì $A-2024B=2035-2024\cdot 0=2035$.
		\end{itemchoice}
	}
\end{ex}

\begin{ex}%[1K6BI-3]
	Nồng độ cồn trong máu (BAC) là chỉ số dùng để đo lượng cồn trong máu của một người. Chẳng hạn, BAC $0{,}02\%$ hay $0{,}2 \mathrm{mg} / \mathrm{ml}$, nghĩa là có $0{,}02 \mathrm{~g}$ cồn trong $100 \mathrm{ml}$ máu. Nếu một người với $\mathrm{BAC}$ bằng $0{,}02 \%$ có nguy cơ bị tai nạn ô tô cao gấp $1{,}4$ lần so với một người không uống rượu, thì nguy cơ tương đối của tai nạn với BAC $0{,}02 \%$ là $1{,}4$. Nghiên cứu y tế gần đây cho thấy rằng nguy cơ tương đối của việc gặp tai nạn khi đang lái ô tô có thể được mô hình hóa bằng một phương trình có dạng
	$$R=\mathrm{e}^{kx},$$
	trong đó $x(\%)$ là nồng độ cồn trong máu và $k$ là một hằng số.\\
	Với giá trị của $k$ được làm tròn đến hàng phần trăm, hãy xét tính đúng sai của các khẳng định sau:
	\choiceTF
	{\True Giá trị của $k$ là $1682,36$}
	{Nguy cơ tương đối của việc gặp tai nạn giao thông là $16,8$ khi nồng độ cồn trong máu là $0{,}17\%$ }
	{\True BAC tương ứng với  nguy cơ tương đối là $100$ là $0,27\%$ (làm tròn đến hàng phần trăm)}
	{Biết rằng một người có nguy cơ tương đối từ $5$ trở lên sẽ không được phép lái xe, thì một người có nồng độ cồn trong máu từ $0,87\%$ trở lên sẽ không được phép lái xe}
	\loigiai{
		\begin{listEX}
			\item Thay $R=1{,}4$ và $x=0{,}02\%$ vào công thức, ta được $1{,}4=\mathrm{e}^{k\cdot \tfrac{0{,}02}{100}}$.\\
			Suy ra $k \approx 1682,36$.
			\item $R=\mathrm{e}^{1682{,}36 \cdot \tfrac{0{,}17}{100}} \approx 17{,}46$.
			\item Thay $R=100$ vào công thức, ta được $100=\mathrm{e}^{1682{,}36 x}$. Suy ra $x \approx 0{,}27 \%$.
			\item Với $R \geq 5$ thì $x \geq 0{,}096 \%$, tức là một người có nồng độ cồn trong máu từ khoảng $0{,}096 \%$ trở lên thì không được lái xe.
		\end{listEX}
	}
\end{ex}

\Closesolutionfile{ans}

\ind{PHẦN III.} \inden{Câu trắc nghiệm trả lời ngắn.}\\
\Opensolutionfile{ans}[ans/1D6-B2-3]

\begin{ex}
	Cho $a,b$ là hai số thực dương, khác $1$. Đặt ${{\log }_{a}}b=2$ , tính giá trị của $P={{\log }_{{{a}^{2}}}}b-{{\log }_{\sqrt{b}}}{{a}^{3}}.$ \\
	\shortans[oly]{$-2$}
	\loigiai
	{
		Ta có:  $P={{\log }_{{{a}^{2}}}}b-{{\log }_{\sqrt{b}}}{{a}^{3}}$. \\
		$=\dfrac{1}{2}{{\log }_{a}}b-6{{\log }_{b}}a$. \\
		$=\dfrac{1}{2}{{\log }_{a}}b-6\dfrac{1}{{{\log }_{a}}b}$. \\
		$=1-3=-2$.
	}
\end{ex}

\begin{ex}
	Cho $\log_a x =-1$ và $\log_a y =4$. Tính giá trị của $P=\log_a (x^2y^3)$.\\
	\shortans[oly]{$10$}
	\loigiai{
		Ta có $P=\log_a (x^2y^3)= 2\log_a x +3 \log_a y =10$. 
	}
\end{ex}

\begin{ex}%[2D2K3-1]%[Đề giữa học kì 1 lớp 12 trường THCS \& THPT Nguyễn Khuyến TP HCM-18]%[WTT2D2-05]%
	Cho $x,y$ là các số dương lớn hơn 1 thỏa mãn $x^2+9y^2=6xy$. Tính $M=\dfrac{1+\log_{12} x+ \log_{12} y}{2 \log_{12} (x+3y)}$.\\
	\shortans[oly]{$1$}
	\loigiai{
		Ta có $x^2+9xy=6xy \Leftrightarrow (x-3y)^2=0 \Leftrightarrow x=3y$.\\
		Khi đó $M=\dfrac{1+2\log_{12} 3y + \log_{12} y}{2\log_{12} 6y} =\dfrac{\log_{12} 12y+ \log_{12} 3y}{\log_{12} 36y^2} =\dfrac{\log_{12} 36y^2}{\log_{12} 36y^2}=1$
	}
\end{ex}

\begin{ex}%[1K6BI-1]
	Biết rằng số chữ số của một số nguyên dương $N$ viết trong hệ thập phân được cho bởi công thức $[\log N]+1$, ở đó $[\log N]$ là phần nguyên của số thực dương $\log N$. Tìm số các chữ số của $2^{2023}$ khi viết trong hệ thập phân. \\
	\shortans[oly]{$609$}
	\loigiai{
		Số chữ số của $2^{2023}$ là $\left[\log 2^{2023}\right]+1=[2023 \cdot \log 2]+1=609$.
	}
\end{ex}

\begin{ex}
	Để đặc trưng đo độ to nhỏ của âm, người ta đưa ra khái niệm mức cường độ âm. Một đơn vị thường dùng để đo mức cường độ âm là đề xin ben (đB). Khi đó, mức cường độ âm $L$ của âm được tính theo công thức $L(I)=10\log\left(\dfrac{I}{I_0} \right)$, trong đó $I$ là cường độ âm tại thời điểm đang xét, $I_0$ là cường độ âm ở ngưỡng nghe ($I_0=10^{-12}\, w/m^2$). Hai cây đàn ghi ta giống nhau, cùng hoà tấu một bản nhạc. Mỗi cây đàn phát ra âm có mức cường độ âm trung bình là 60 đB. Hỏi mức cường độ âm tổng cộng do hai chiếc đàn cùng phát ra là bao nhiêu đề xin ben (đB)?\\
	\shortans[oly]{$63$}
	\loigiai{
	}
\end{ex}

\begin{ex}%[1D4B2-4]
	Để tính độ tuổi của mẫu vật bằng gỗ, người ta đo độ phóng xạ của $_{~6}^{14}\mathrm{C}$  có trong mẫu vật tại thời điểm $t$ (năm) (so với thời điểm ban đầu $t=0$ ), sau đó sử dụng công thức tính độ phóng xạ $H=H_0 \mathrm{e}^{-\lambda t}$ (đơn vị là Becquerel, kí hiệu $\mathrm{Bq}$ ) với $H_0$ là độ phóng xạ ban đầu (tại thời điểm $t=0$); $\lambda=\dfrac{\ln 2}{T}$ là hằng số phóng xạ, $T=5730$ (năm) (Nguồn: Vật lí 12 Nâng cao, NXBGD Việt Nam, 2014). Khảo sát một mẫu gỗ cổ, các nhà khoa học đo được độ phóng xạ là $0{,}215$ Bq. Biết độ phóng xạ của mẫu gỗ tươi cùng loại là $0{,}250$ Bq. Xác định độ tuổi (tính bằng năm) của mẫu gỗ cổ đó (làm tròn kết quả đến hàng đơn vị).\\
	\shortans[oly]{$1 247$}
	\loigiai{
		Ta có độ phóng xạ của mẫu gỗ cổ ở thời điểm ban đầu ($t=0$) là $H_0=0{,}250$ Bq.\\
		Độ phóng xạ của mẫu gỗ cổ ở thời điểm khảo sát là $H=0{,}215$ Bq.\\
		Ta có $H=H_0\mathrm{e}^{-\lambda t}\Leftrightarrow 0{,}215=0{,}250\cdot e^{-\lambda t}=0{,}250\cdot 2^{-\tfrac{t}{T}}\Leftrightarrow \dfrac{t}{T}=-\log_2 \left( \dfrac{0{,}215}{0{,}250}\right)\Rightarrow t\approx 1 247$ (năm).\\
		Vậy tuổi của mẫu gỗ cổ đó là $1247$ năm.
	}
\end{ex}

\Closesolutionfile{ans}

